% chapters/04-server-actions.tex -- group B (S2): Server Actions와 데이터 읽기. Items B1-B4.
% ONLY sources: perf.sources.md S2, perf.audit.md §0-§2 group B. Items are copied.
\chapter{Server Actions와 데이터 읽기}\label{ch:actions}

% src: perf.audit.md §2 group B, §3 (B: PASS 2, PARTIAL 1, FAIL 1)
\paragraph{결과.}
묶음 B는 4개 중 2개가 통과다. 부분 통과 1개, 실패 1개다.
쓰이는 action 4개는 모두 변경(PATCH, DELETE)이다. 읽기 action은 정의만 남은 죽은 코드 1개뿐이다.
남은 문제는 무한 스크롤의 load-more가 edge cache 밖이라는 점이다.

\begin{inbrief}
  \item[Why] Server Action은 줄을 서므로 읽기에 쓰면 시간이 더해진다(S2).
  \item[When] 한 화면이 읽기를 여러 개 하거나 첫 렌더에서 읽을 때다.
  \item[Where] \texttt{app/actions/}, TanStack Query hook, \texttt{app/api/}의 Route Handler다.
  \item[How] 변경만 action으로 한다. 읽기는 Server Component나 Route Handler로 한다.
\end{inbrief}

\section{원리: action은 줄을 선다}\label{sec:ac-principle}

\paragraph{문제와 원인.}
\begin{itemize}[nosep]
  \item Server Action은 client마다 하나씩 줄을 선다.
  \item 1초짜리 action 셋을 \texttt{Promise.all}로 불러도 약 3초다.
  \item 첫 렌더 중에는 부를 수 없다. \texttt{queryFn}으로 쓰면 \texttt{useSuspenseQuery}가 실패한다.
\end{itemize}
% src: perf.sources.md S2 bullets 1, 3

\paragraph{무엇을 알려고 쓰나.}
같은 읽기 셋을 action으로 하면 왜 3초이고, Server Component나 Route Handler로 하면 왜 줄어드는가.
\Cref{fig:ac-flow}가 두 흐름을 나란히 보인다.

\begin{figure}[H]
  \centering
  \resizebox{\linewidth}{!}{%
  \begin{tikzpicture}[font=\scriptsize]
    \node[font=\small\bfseries, anchor=west] at (-1.0,1.2) {읽기를 Server Action으로 (\texttt{Promise.all}로 불러도)};
    \node[badbox, minimum width=2.4cm] (b1) at (0,0) {action 1\\1초};
    \node[badbox, minimum width=2.4cm] (b2) at (3.4,0) {action 2\\1초};
    \node[badbox, minimum width=2.4cm] (b3) at (6.8,0) {action 3\\1초};
    \node[flowbox, minimum width=2.4cm] (b4) at (10.2,0) {합계 약 3초};
    \draw[flowarrow] (b1) -- (b2);
    \draw[flowarrow] (b2) -- (b3);
    \draw[flowarrow] (b3) -- (b4);

    \node[font=\small\bfseries, anchor=west] at (-1.0,-1.7) {읽기를 Server Component나 Route Handler로};
    \node[goodbox, minimum width=2.4cm] (a1) at (0,-3.0) {읽기 1};
    \node[goodbox, minimum width=2.4cm] (a2) at (3.4,-3.0) {읽기 2};
    \node[goodbox, minimum width=2.4cm] (a3) at (6.8,-3.0) {읽기 3};
    \node[flownew, minimum width=2.4cm] (a4) at (10.2,-3.0) {줄 서지 않음};
    \draw[flowarrow] (a1) -- (a4);
    \draw[flowarrow] (a2) -- (a4);
    \draw[flowarrow] (a3) -- (a4);
  \end{tikzpicture}}
  \par\smallskip{\centering\scriptsize \flowkey{badbox} 줄 서는 action \quad \flowkey{goodbox} 줄 서지 않는 읽기 \quad \flowkey{flownew} 바뀐 단계\par}
  \caption{읽기 셋의 두 방식. 위는 action이 줄을 서서 약 3초다.}
  \label{fig:ac-flow}
  % src: perf.sources.md S2 (three 1 s actions via Promise.all ≈ 3 s; reads through Server Components / Route Handlers). The "not queued" box is the book's drawing of the S2 design rule.
\end{figure}

\paragraph{처방.}
\begin{enumerate}[nosep]
  \item Server Action은 변경(mutation) 전용으로 쓴다.
  \item 읽기는 Server Component나 Route Handler로 한다.
  \item backend 호출은 Next server를 거치는 BFF(backend for frontend)로 한다.
  \item BFF는 Next server와 같은 곳에 있어 왕복이 약 3 ms다. 브라우저에서 backend로 가면 왕복마다 약 100 ms다.
\end{enumerate}
% src: perf.sources.md S2 bullets 4-5

\section{체크리스트 B1--B4}\label{sec:ac-checklist}

\begin{checkitem}{\ckdone}{B1}{읽기는 Server Action으로 하지 않는다. action을 병렬로 부르지 않는다}{PASS\ckdot{}영향 없음}
  % src: perf.audit.md B1
  \item[What] 읽기는 RSC나 Route Handler로 한다. action은 변경 전용이다.
  \item[Why] S2. action은 줄을 선다. 첫 렌더에 부를 수 없다.
  \item[How] \code{rg -l "'use server'"}가 2 파일을 낸다. 쓰이는 action 4개는 모두 PATCH나 DELETE다(\code{app/actions/postPublish.action.ts:30,45,58,74}). \code{rg -n "Promise\.all" app components hooks apis}는 0건이다.
  \item[When] 유지.
\end{checkitem}

\begin{checkitem}{\ckopen}{B2}{죽은 읽기 action이 없다}{FAIL\ckdot{}영향 L (1건)}
  % src: perf.audit.md B2
  \item[What] 쓰이지 않는 \code{'use server'} 읽기 함수가 없다.
  \item[Why] S2. 읽기 action은 잘못된 본보기다. 다시 쓰이면 순차 실행이 된다.
  \item[How] \code{rg -n loadMorePosts}가 1건을 낸다. 정의뿐이다(\code{app/actions/post.action.ts:6}). 처방: \code{x_post.action.ts}로 이름을 바꾼다(soft-delete).
  \item[When] 지금.
\end{checkitem}

\begin{checkitem}{\ckdone}{B3}{action을 \texttt{queryFn}으로 쓰지 않는다}{PASS\ckdot{}영향 없음}
  % src: perf.audit.md B3
  \item[What] TanStack \code{queryFn}은 fetch 함수다. action은 \code{mutationFn}에만 쓴다.
  \item[Why] S2. \code{queryFn}의 action은 Suspense query에서 깨진다.
  \item[How] action을 import하는 hook이 4개다. \code{hooks/useDeletePost.query.client.ts:25}, \code{useDeleteAdminPostAction.query.client.ts:19}, \code{useSetAdminPostStatus.query.client.ts:27}은 모두 mutation이다. \code{useEditorHandlers.hook.ts:182}는 event handler다.
  \item[When] 유지.
\end{checkitem}

\begin{checkitem}{\ckopen}{B4}{load-more BFF가 edge cache와 rate limit을 고려한다}{PARTIAL\ckdot{}영향 M (edge 상태는 확인 못함, \Cref{sec:pl-open})}
  % src: perf.audit.md B4
  \item[What] 무한 스크롤 읽기가 cache되거나, 요청 수 한도에 걸리지 않는다.
  \item[Why] S2는 Next server를 거치는 BFF를 권한다. 하지만 여기서 \code{/api/}는 edge cache 밖이다. 스크롤 한 page가 VM 요청 하나이고 50 / 10 s에 센다(\decision{68}, \decision{70}의 429).
  \item[How] client \code{components/blog/RecentPostsList.client.tsx:62}와 \code{TagPostsList.client.tsx:60}이 \code{/api/posts/load-more}를 부른다. Cloudflare 규칙은 \code{deploy/traffic/README.md:76}의 \code{not starts_with(..., "/api/")}다. \code{app/api/posts/load-more/route.ts:49,56}은 \code{Cache-Control}을 붙이지 않는다. 더해 이 route는 token 없이 읽는다(\code{:52}). 첫 page(\code{getFeed}, token 있음)와 다르다. 소유자의 draft가 2쪽부터 빠진다. 처방 후보: (a) 익명 응답에 \code{Cache-Control: public, s-maxage=60}과 Cloudflare 규칙을 더한다. (b) 브라우저가 \code{blog-api.} GET을 직접 부른다(이미 edge \cfs{HIT}).
  \item[When] 나중(traffic 설계와 함께). (a)와 (b)의 선택은 소유자.
\end{checkitem}
